% S1 Text, part 1: datasets, models, decoding, scoring, detector and tracer results.

This text gives the settings needed to repeat the training, decoding and scoring of the three published models.
The models are the detector (StarSeg), the field network (Hydra) and the tracer (Noodle Tracer).
Together they trace every arm of the sunflower sea star (\textit{Pycnopodia helianthoides}) in images and video.
All results come from one training campaign on two frozen datasets, and the published models are seed 0 of that campaign.
Precision, recall and F1 are proportions between 0 and 1.
Differences between tracer configurations are given in percentage points.
Distances are in pixels (px) of the 512 $\times$ 512 crop that the field network reads, unless stated otherwise.
An acquisition group is the unit by which the arm-annotation dataset was split.
A capture group is the unit by which the detection dataset was split.
Neither unit is an animal.

\section*{Datasets}

\subsection*{Arm-annotation dataset}

The arm-annotation dataset holds 5,093 square RGB crops, each showing one sunflower sea star (Table~\ref{tab:armdata}).
Crops are stored at native resolution.
The median side is 519~px (quartiles 174 and 1,264~px; range 71 to 13,000~px).
The models read each crop resized to 512 $\times$ 512~px, and scoring uses the same canvas.

Annotation was done in CVAT~\cite{cvat2026computer}.
Each arm is one polyline ordered from the disc to the tip, and the order defines progress along the arm.
The dataset has 57,055 arm polylines, with a median of 10 per image (range 3 to 28).
Points mark the mouth (399 points), the madreporite (1,254) and the anus (4,693) on the images where they were annotated.
\gap{state the rule for marking landmarks, the number of annotators and any agreement check}
Neighbouring stars that were not annotated are covered by ignore polygons (515 polygons in 329 images).
Pixels inside these polygons do not enter the loss of the field network.

The crops come from seven sources (Table~\ref{tab:armdata}).
Two sources are crops of single stars from images of the two annotation projects of the detection dataset (\texttt{starseg} and \texttt{starseg\_lab}).
Two smaller sources are named \texttt{arm\_pit} and \texttt{axis\_detector}. \gap{describe the \texttt{arm\_pit} and \texttt{axis\_detector} image sets}
Sleepy Stars, StarDome and StarTrack are laboratory camera rigs. \gap{cameras, lenses, frame rates and lighting of each source}
For Sleepy Stars a video frame was kept only if a detector found exactly one sunflower sea star with a confidence of at least 0.90.
The detector also had to find the star in both adjacent frames, 1~s apart, with a confidence of at least 0.75.
The crop is centred on the bounds of the detected polygon, its side is 1.4 times the larger polygon dimension, and it had to fit inside the frame.
\gap{crop rule for the crops from detector images, StarDome and StarTrack}

The split unit is the acquisition group.
A group is the set of images that share an exact source key, a session key or identical bytes.
Whole groups were allocated to one split with seed 20261026, which gave 394 training, 180 validation and 211 test groups.
No group appears in two splits.
Groups are recording sessions, image clusters or annotation tasks.
They are not animals, and the same animal can appear in more than one split. \gap{number of animals in each source, if known}
StarDome and StarTrack each contribute 70 training and 30 test images and no validation image.
The StarTrack test images are all images of camera C and form one group; cameras A, B and D are in training.
The StarDome test groups are single images.
For both sources the independence of test and training animals was not verified.

Before the split was frozen we searched for near-duplicate images across splits.
We found 17 cross-split pairs with an identical perceptual hash.
Their 12 connected components touched 28 groups, and we removed 149 images (140 training, 9 test).
The split went from 4,202, 478 and 562 images to 4,062, 478 and 553.

Of the 553 test images, 493 (180 groups) come from the four sources that also have validation images.
We call them the development test set, because they were scored several times while the method was developed.
The development test set is the primary cohort for the comparisons between tracer configurations.
Three further caveats apply to it.
About 47 test images have a near-identical training image in another group, because the removal step used identical hashes only.
Two audits counted 217 and 218 test crops that were cut from images in the training split of the detection dataset; both counts predate the final version of that dataset.
The acquisition groups do not separate animals.

\begin{table}[!htbp]
\centering
\caption{\textbf{Arm-annotation dataset by source and split.} Cells give images, with acquisition groups in parentheses.}
\label{tab:armdata}
{\small
\begin{tabular}{|l|r|r|r|r|r|}
\hline
\textbf{Source} & \textbf{Training} & \textbf{Validation} & \textbf{Test} & \textbf{All images} & \textbf{All arms} \\ \thickhline
Detector images, \texttt{starseg} & 2,366 (158) & 255 (99) & 248 (99) & 2,869 & 31,304 \\ \hline
Detector images, \texttt{starseg\_lab} & 732 (124) & 138 (73) & 138 (72) & 1,008 & 8,950 \\ \hline
\texttt{arm\_pit} & 394 (19) & 51 (5) & 44 (4) & 489 & 8,094 \\ \hline
\texttt{axis\_detector} & 27 (1) & 0 & 0 & 27 & 400 \\ \hline
Sleepy Stars & 403 (19) & 34 (3) & 63 (5) & 500 & 5,818 \\ \hline
Subtotal, development sources & 3,922 (321) & 478 (180) & 493 (180) & 4,893 & 54,566 \\ \thickhline
StarDome & 70 (70) & 0 & 30 (30) & 100 & 1,387 \\ \hline
StarTrack & 70 (3) & 0 & 30 (1) & 100 & 1,102 \\ \thickhline
All & 4,062 (394) & 478 (180) & 553 (211) & 5,093 & 57,055 \\ \hline
Arm polylines & 46,569 & 4,801 & 5,685 & & 57,055 \\ \hline
Arms per image, median (range) & 11 (3 to 28) & 9 (4 to 23) & 10 (3 to 22) & & \\ \hline
\end{tabular}
}
\begin{flushleft}\small
The test images of the four development sources that have test images form the development test set (493 images, 180 groups).
The subtotal row sums the rows above it.
Source names in typewriter type are the labels used in the dataset files.
\end{flushleft}
\end{table}

\subsection*{Detection dataset}

The detection dataset holds 7,978 whole images with 26,394 instance masks, annotated in CVAT.
There are 37 classes: the sunflower sea star (8,432 masks), 27 other sea star classes (26 species and one class for unknown species), six classes of live prey and three of dead prey.
We trained with four label schemes on the same images.
Scheme F keeps all 37 classes.
Scheme C has four classes (sunflower sea star, other sea star, live prey, dead prey).
Scheme B has two classes (sunflower sea star, other object).
Scheme P labels sunflower sea stars only and leaves other objects unlabelled; the images are kept.
Images were resized so that the longest side is at most 1,024~px and stored as JPEG at quality 95.
Scoring uses the original annotation shapes in source pixel coordinates.

Each image received a capture group by a rule specific to its source.
Examples are one group per web photograph, per capture day, per source video, per camera and trial, or per rig.
Groups were then merged across exact duplicate files (169 links) and across near duplicates (2,105 links).
Near duplicates were defined by a 256-bit difference hash with a Hamming distance of at most 8.
Merging reduced 3,607 capture groups to 3,494.
Whole merged groups were allocated to one split with seed 20261023 (Table~\ref{tab:detdata}).
The split has 5,600 training, 1,129 validation and 1,249 test images in 2,426, 548 and 520 groups.
Sunflower sea star masks number 5,806, 1,213 and 1,413.
Images without a sunflower sea star number 2,771, 642 and 550.
Of these, 84 training images and 8 test images carry no annotation at all, and each was checked by eye.
The validation split has no annotation-free image, so false alarms on empty scenes cannot be measured there.

We group the images in two ways.
S1 Fig uses five image domains: laboratory, field survey video, captive tanks outside the laboratory, centred field photographs and public web photographs.
The evaluation uses a coarser view of the test split: Lab, Field and Web.
Field is the union of the three middle domains (395 test images, 8 groups, 427 stars).
We split Field further into PWS (Prince William Sound, Alaska) and field without PWS, and we never pool the two.

PWS is a set of 878 diver photographs, each framed on one sunflower sea star.
They were taken with an Olympus TG-6 camera on dive surveys during a two-night cruise in the sound, on three days from 21 to 23 September 2023.
The first day is the test split (243 images, one capture group, 38 individual stars).
The other two days are in training (635 images with 646 masks).
Training and test images therefore come from one cruise, and the PWS score is a within-cruise result.
No individual star is in both splits by its field identifier, but identifiers were not matched across days.
Of the PWS images, 523 were drafted from box proposals with the SAM2 segmentation model and then annotated and checked by the author in CVAT.

Two further caveats apply.
Of the 7,978 images, 7,881 were used in earlier fits of the detector, so the test split is held out for this campaign only.
The origin of the web photographs is not recorded. \gap{origin and licences of the web photographs and contributed photographs}

\begin{table}[!htbp]
\centering
\caption{\textbf{Detection dataset by image domain and split.} Stars are sunflower sea star masks; groups are capture groups after merging of duplicates.}
\label{tab:detdata}
{\footnotesize
\setlength{\tabcolsep}{3.5pt}
\begin{tabular}{|l|r|r|r|r|r|r|r|r|r|}
\hline
 & \multicolumn{3}{c|}{\textbf{Training}} & \multicolumn{3}{c|}{\textbf{Validation}} & \multicolumn{3}{c|}{\textbf{Test}} \\ \hline
\textbf{Domain} & \textbf{Images} & \textbf{Stars} & \textbf{Groups} & \textbf{Images} & \textbf{Stars} & \textbf{Groups} & \textbf{Images} & \textbf{Stars} & \textbf{Groups} \\ \thickhline
\multicolumn{10}{|l|}{\textbf{Five image domains}} \\ \hline
Laboratory & 1,951 & 4,721 & 47 & 341 & 1,065 & 17 & 352 & 958 & 25 \\ \hline
Field survey video & 339 & 159 & 14 & 64 & 48 & 4 & 32 & 30 & 2 \\ \hline
Captive tanks outside the laboratory & 23 & 45 & 1 & 72 & 72 & 1 & 55 & 135 & 2 \\ \hline
Centred field photographs & 859 & 752 & 5 & 112 & 0 & 1 & 308 & 262 & 4 \\ \hline
Public web photographs & 2,428 & 129 & 2,359 & 540 & 28 & 525 & 502 & 28 & 487 \\ \hline
All & 5,600 & 5,806 & 2,426 & 1,129 & 1,213 & 548 & 1,249 & 1,413 & 520 \\ \thickhline
\multicolumn{10}{|l|}{\textbf{Evaluation view}} \\ \hline
Lab & 1,951 & 4,721 & 47 & 341 & 1,065 & 17 & 352 & 958 & 25 \\ \hline
Field without PWS & 586 & 310 & 18 & 248 & 120 & 6 & 152 & 184 & 7 \\ \hline
PWS & 635 & 646 & 2 & 0 & 0 & 0 & 243 & 243 & 1 \\ \hline
Web & 2,428 & 129 & 2,359 & 540 & 28 & 525 & 502 & 28 & 487 \\ \hline
\end{tabular}
}
\begin{flushleft}\small
Lab equals the laboratory domain and Web equals the public web photographs.
Field is the union of survey video, captive tanks and centred field photographs.
PWS is part of the centred field photographs.
The training and validation counts of the row ``Field without PWS'' are differences between the rows above and the PWS row.
All instance masks of any class: 18,345 training, 3,966 validation, 4,083 test.
\end{flushleft}
\end{table}

\section*{Models and training}

Table~\ref{tab:settings} lists the training settings of the three models side by side.

\subsection*{Detector}

The detector is a YOLO26-L segmentation model with an end-to-end head, as implemented in Ultralytics~\cite{jocher2026ultralytics}.
It has 31,375,206 parameters and was initialized from weights pretrained on COCO~\cite{lin2014microsoft}.
We trained 60 models with one recipe: five model sizes (N, S, M, L, X) by four label schemes by three seeds.
Training used AdamW~\cite{loshchilov2019decoupled} with a cosine schedule~\cite{loshchilov2017sgdr} and mixed precision~\cite{micikevicius2018mixed}.
The loss gains were 7.5 (box), 0.5 (class) and 1.5 (distribution focal loss).
Every model was scored from its terminal checkpoint, with no selection of a best epoch.

We registered the selection rule before any test prediction was made.
On 11 November 2026 the rule was registered at 14:58:34 UTC, the selection was locked at 15:57:31 and the first test prediction started at 16:07:50.
The test split was read once for all 60 models.
The rule selects the configuration with the highest mean validation mask AP for the sunflower sea star class over three seeds.
Mask AP is the COCO average precision over intersection-over-union thresholds from 0.50 to 0.95~\cite{lin2014microsoft}, computed with pycocotools 2.0.11.
Ties go to the configuration with fewer parameters and then to the scheme that comes first in alphabetical order.
The rule selected size L with scheme B (mean 0.6340, SD 0.0025 over seeds; runner-up X with scheme C, 0.6330).
The published detector is seed 0 of that configuration.
For mask AP, predictions were made at 32-bit precision with an image size of 1,024, a confidence of 0.001 and at most 300 detections per image.

At inference the detector runs at an image size of 1,024 with a confidence of 0.25 and an intersection-over-union setting of 0.7.
Only detections of the sunflower sea star class are kept.
Each detection with box width $w$ and height $h$ gives one square crop centred on the box centre, with side
\begin{equation}
s = \max\!\left(16,\ \left\lceil 2.0\,\max(w,h)\right\rceil\right)\ \text{px}.
\label{eq:s1crop}
\end{equation}
Parts of the square outside the frame are filled with zeros.
The crop is resized to 512 $\times$ 512~px with bilinear interpolation.
A point at crop coordinates $(x, y)$ maps back to the frame as $(x, y)\,s/512 + (x_0, y_0)$, where $(x_0, y_0)$ is the top left corner of the crop.

\subsection*{Field network}

The field network maps a crop to 18 maps of the same size (Table~\ref{tab:fields}).
The backbone is a ConvNeXt-Base~\cite{liu2022convnet} with DINOv3 pretraining~\cite{simeoni2025dinov3}.
It has 87,566,464 parameters and stage outputs of 128, 256, 512 and 1,024 channels, and all of it is trained.
The input is normalized inside the backbone with the ImageNet mean and standard deviation~\cite{deng2009imagenet}.
A feature pyramid neck~\cite{lin2017feature} with 1,673,216 parameters projects each stage to 256 channels with a 1 $\times$ 1 convolution.
It adds the stages from the top down after bilinear upsampling.
It then refines the sum at the resolution of the first stage with two blocks of 3 $\times$ 3 convolution, group normalization~\cite{wu2018group} and SiLU~\cite{elfwing2018sigmoid}.
Five heads share one design.
A 1 $\times$ 1 convolution from 256 to 128 channels is followed by a residual block of two 3 $\times$ 3 convolutions with group normalization, then SiLU and a 1 $\times$ 1 convolution to the outputs.
The heads have 328,578 (support and centreline), 328,707 (geometry), 328,707 (landmark), 328,578 (endpoint) and 329,352 (embedding) parameters.
Their outputs are concatenated and upsampled bilinearly to 512 $\times$ 512.
The network has 90,883,602 parameters in total.

Before the targets are drawn, each arm polyline is smoothed with two Chaikin iterations~\cite{chaikin1974algorithm} and resampled at 1~px.
The loss is a weighted sum of six terms,
\begin{equation}
L = 1.0\,L_{\mathrm{support}} + 3.0\,L_{\mathrm{centreline}} + 1.0\,L_{\mathrm{geometry}} + 10.0\,L_{\mathrm{landmark}} + 10.0\,L_{\mathrm{endpoint}} + w_e\,L_{\mathrm{embedding}},
\label{eq:s1loss}
\end{equation}
with $w_e = 0.5\,\min(1, \mathrm{epoch}/4)$.
The support term is a binary cross-entropy with positive weight $\mathrm{clamp}(n_{\mathrm{neg}}/n_{\mathrm{pos}}, 1, 100)$, plus a soft Dice loss~\cite{milletari2016vnet}, plus 2.0 times an area penalty.
The area penalty is $\mathrm{relu}(\bar{p} - 1.15\,\bar{y})^2/\bar{y}$, where $\bar{p}$ and $\bar{y}$ are the mean predicted and target support.
The centreline term is a focal binary cross-entropy~\cite{lin2017focal} ($\alpha = 0.75$, $\gamma = 2.0$) averaged over the pixels of the target support.
The geometry term is an L1 loss on the tangent and on progress, again on the pixels of the target support.
The landmark and endpoint terms are mean squared errors with pixel weight $1 + 20 \times \mathrm{target}$, normalized by the sum of the weights.
The embedding term is a discriminative loss~\cite{debrabandere2017semantic}.
For each image it pulls the embedding $e$ of every arm pixel towards the mean $\mu$ of its arm with $\max(0, \lVert e - \mu\rVert - 0.5)^2$.
It pushes the means of every pair of arms apart with $\max(0, 2 \times 1.5 - \lVert \mu_a - \mu_b\rVert)^2$, and adds $0.001$ times the mean norm of the arm means.
Pixels covered by more than one arm are left out of the embedding term.

The backbone learning rate is $10^{-5}$ in the deepest stage and falls by a factor of 0.7 per stage towards the stem ($10^{-5}$, $7 \times 10^{-6}$, $4.9 \times 10^{-6}$, $3.43 \times 10^{-6}$).
The output biases of the support and centreline, landmark and endpoint heads start at $-4.0$.
The published weights are an exponential moving average of the trained weights~\cite{polyak1992acceleration} with decay 0.999, taken at the last step.
No gradient clipping was used.
Geometric augmentation is a horizontal flip (probability 0.5), a vertical flip (0.15) and a rotation drawn uniformly from $-180$ to 180 degrees.
Photometric augmentation scales brightness and contrast by a factor from U(0.82, 1.18) and saturation by U(0.80, 1.20), and shifts hue by up to 0.03 of the hue circle.
It also converts to greyscale with probability 0.08, applies a gamma from U(0.85, 1.15), blurs with probability 0.15 (radius up to 1~px) and adds Gaussian noise with SD 0.015.
A simulation of a low-quality camera is applied with probability 0.4 to images whose shorter native side is at least 384~px.
It blurs the image (probability 0.7, radius 0.5 to 2.5~px) and reduces it so that the longer side is between 96 and 256~px.
It then adds noise with SD between 0.005 and 0.03, compresses the image as JPEG at quality 40 to 80 (probability 0.5) and enlarges it back.
Last, it scales contrast by 0.45 to 0.85 (probability 0.6) and brightness by 0.55 to 0.95 (0.6), and converts the image to greyscale (0.5).
We trained three field networks (seeds 0, 1 and 2) with this recipe; the published one is seed 0.

\begin{table}[!htbp]
\centering
\caption{\textbf{Training settings of the three published models.}}
\label{tab:settings}
{\footnotesize
\setlength{\tabcolsep}{3.5pt}
\begin{tabular}{|>{\raggedright\arraybackslash}p{0.85in}|>{\raggedright\arraybackslash}p{1.6in}|>{\raggedright\arraybackslash}p{1.85in}|>{\raggedright\arraybackslash}p{1.65in}|}
\hline
\textbf{Setting} & \textbf{Detector} & \textbf{Field network} & \textbf{Tracer} \\ \thickhline
Training data & 5,600 images, label scheme B & 4,062 crops & 4,062 crops with 46,569 arms; maps of the seed 0 field network \\ \hline
Validation data & 1,129 images & 478 crops & 478 crops with 4,801 arms \\ \hline
Input & image size 1,024 & 512 $\times$ 512 crop & 33 $\times$ 33 patches of the maps and of image features \\ \hline
Initialization & YOLO26-L segmentation weights pretrained on COCO & backbone from DINOv3 pretraining; neck and heads random & random; image encoder from ImageNet pretraining, frozen \\ \hline
Optimizer & AdamW, momentum 0.9, weight decay 0.0005 & AdamW, weight decay $10^{-4}$ & AdamW, weight decay $10^{-4}$, gradient norm clipped at 5.0 \\ \hline
Learning rate & 0.000244, final factor 0.01 & neck and heads $10^{-4}$; backbone $10^{-5}$ with factor 0.7 per stage & $3 \times 10^{-4}$ \\ \hline
Schedule & warm-up of 3 epochs, then cosine & linear warm-up over 800 steps from a factor of 0.02, then cosine to a factor of 0.05 & per cycle: factor 0.5 in cycle 1, 1.0 in cycle 2, then cosine to 0.1 in cycle 30 \\ \hline
Length & 50 epochs, no early stopping & 98,250 steps (194 epochs of 508 steps) & 30 cycles of 4,000 steps (120,000 steps) \\ \hline
Batch & 4 images, nominal batch size 64 & 8 crops & 1 image, all its seeds \\ \hline
Augmentation & hue 0.015, saturation 0.7, value 0.4; translation 0.1; scale 0.5; horizontal flip 0.5; mosaic 1.0, off in the last 10 epochs; RandAugment; random erasing 0.4 & flips, rotation by any angle, photometric changes, simulation of a low-quality camera (see text) & new draw for every image in each cycle: flips on each axis (0.5), rotation by a multiple of 90 degrees, brightness, contrast and saturation factors from U(0.85, 1.15) \\ \hline
Mixed precision & on & on; embedding loss at 32-bit & on \\ \hline
Checkpoint rule & terminal checkpoint & terminal checkpoint, moving average of weights (decay 0.999) & terminal checkpoint after cycle 30 \\ \hline
Hardware & NVIDIA RTX A6000 & NVIDIA A40, rented & workstation GPU with 32,607~MiB (RTX 5090, inferred from the memory size) \\ \hline
Wall time & about 6.43~h & 85,097~s (about 23.6~h) & 30,616~s (about 8.5~h), including regeneration of the maps \\ \hline
Peak GPU memory & about 13.45~GiB & 12,310~MiB & not recorded for this job alone \\ \hline
Software & Python 3.12.3, PyTorch 2.8.0, Ultralytics 8.4.51 & PyTorch 2.6.0 as pinned; the version at run time was not recorded & Python 3.12.3, PyTorch 2.8.0 \\ \hline
Parameters & 31,375,206 & 90,883,602 & 194,463, of which 183,975 trained \\ \hline
\end{tabular}
}
\begin{flushleft}\small
All three models are seed 0 of their configuration.
Detector augmentation values are the Ultralytics settings recorded for the run.
PyTorch is cited as~\cite{paszke2019pytorch}.
\end{flushleft}
\end{table}

\begin{table}[!htbp]
\centering
\caption{\textbf{The 18 output maps of the field network.} Loss weights are those of Eq~\ref{eq:s1loss}.}
\label{tab:fields}
{\footnotesize
\setlength{\tabcolsep}{3.5pt}
\begin{tabular}{|l|>{\raggedright\arraybackslash}p{0.9in}|>{\raggedright\arraybackslash}p{0.75in}|l|>{\raggedright\arraybackslash}p{2.0in}|>{\raggedright\arraybackslash}p{1.25in}|}
\hline
\textbf{Channel} & \textbf{Map} & \textbf{Head} & \textbf{Output} & \textbf{Training target} & \textbf{Loss (weight)} \\ \thickhline
0 & Arm support & Support and centreline & sigmoid & union of capsules 8~px wide around each arm polyline & weighted cross-entropy, soft Dice and area penalty (1.0) \\ \hline
1 & Arm centreline & Support and centreline & sigmoid & $\mathrm{clip}(1 - (d/1.5)^2, 0, 1)^2$, with $d$ the distance in px to the centreline; zero outside the support & focal cross-entropy on support pixels (3.0) \\ \hline
2, 3 & Tangent, $x$ and $y$ & Geometry & tanh & unit tangent pointing from base to tip, on the support & L1 on support pixels (1.0, with progress) \\ \hline
4 & Progress & Geometry & sigmoid & arc length along the arm, from 0 at the base to 1 at the tip & L1 on support pixels (1.0, with tangent) \\ \hline
5, 6, 7 & Mouth, madreporite, anus & Landmark & sigmoid & Gaussian with SD 4.0~px at each annotated point & peak-weighted squared error (10.0) \\ \hline
8 & Arm base & Endpoint & sigmoid & Gaussian with SD 1.5~px at the first point of each polyline & peak-weighted squared error (10.0, with tip) \\ \hline
9 & Arm tip & Endpoint & sigmoid & Gaussian with SD 2.0~px at the last point of each polyline & peak-weighted squared error (10.0, with base) \\ \hline
10 to 17 & Embedding, 8 dimensions & Embedding & linear & none; pixels are labelled by arm & discriminative loss (0.5 after a ramp over 4 epochs) \\ \hline
\end{tabular}
}
\begin{flushleft}\small
Tangent and progress values are averaged where arm segments overlap.
\end{flushleft}
\end{table}

\subsection*{Tracer}

The tracer walks along each arm from the tip towards the disc, using the maps of the field network and the crop.
Seeds are local maxima of the tip map with a score of at least 0.2, after non-maximum suppression with a radius of 4~px.
The seed embedding $e$ is the mean embedding over pixels with support above 0.2 within 2~px of the seed.

At each step the tracer reads a 33 $\times$ 33 patch of the 18 maps, centred on the current position.
The maps are taken before their output activations and are standardized per channel.
The mean and SD for this come from the 3,659 training crops that are not from Sleepy Stars, without augmentation.
Three channels are added: the similarity $\exp(-0.5\,\lVert e(p) - e\rVert^2)$ of each pixel embedding $e(p)$ to the seed embedding, and two coordinate ramps from $-1$ to 1.
A stem of two strided convolutions (21 to 32 channels, 5 $\times$ 5; 32 to 64 channels, 3 $\times$ 3; stride 2 each) with group normalization and SiLU reduces the 21 channels to a 9 $\times$ 9 grid.
Image features come from the first layers of a MobileNetV3-Small~\cite{howard2019searching} pretrained on ImageNet~\cite{deng2009imagenet} and kept frozen.
Two feature maps (16 and 24 channels) are computed once per crop, sampled as 9 $\times$ 9 patches at the current position and projected to 16 channels each.
A 3 $\times$ 3 convolution fuses the 96 channels to 64, giving features $F_t$ at 81 positions.

An attention step~\cite{vaswani2017attention} pools the 81 positions into a context vector $c_t$.
A gated recurrent unit~\cite{cho2014learning} with 128 hidden units updates the state $h_t$ from the context, the previous displacement $\Delta_{t-1}$ and the seed embedding:
\begin{equation}
\begin{aligned}
c_t &= \sum_{k=1}^{81} \operatorname{softmax}_k\!\left(\frac{K(F_t)_k^{\top}\, q(h_{t-1},\Delta_{t-1},e)}{\sqrt{64}}\right) F_{t,k}, \\
h_t &= \operatorname{GRU}\!\left([\,c_t,\ \Delta_{t-1},\ e\,],\ h_{t-1}\right), \\
\Delta_t &= \frac{W_d h_t + b_d}{\lVert W_d h_t + b_d\rVert_2}\,\bigl(0.5 + 5.5\,\sigma(w_s^{\top} h_t + b_s)\bigr), \qquad x_{t+1} = x_t + \Delta_t .
\end{aligned}
\label{eq:s1step}
\end{equation}
Here $K$ is a 1 $\times$ 1 convolution, $q$ is a linear layer, $\sigma$ is the logistic function, and $h_0$ and $\Delta_0$ are zero.
The displacement is between 0.5 and 6~px per step.
A third linear head gives the probabilities of four states: invalid seed, continue, terminate at base and terminate lost.
The tracer has 194,463 parameters, of which 10,488 belong to the frozen image encoder.
The field network and the tracer together have 91,078,065 parameters.

Training uses the maps of the frozen seed 0 field network.
It runs for 30 cycles of 4,000 steps.
Before each cycle the maps are computed again for one new augmented version of every training crop (Table~\ref{tab:settings}; augmentation seed 20261026).
One training step uses all seeds of one image.
Tip peaks are matched one to one to annotated tips within 12~px.
Unmatched peaks within 20~px of an annotated tip are ignored, and peaks farther away serve as examples of invalid seeds.
Start positions are jittered with a Gaussian of SD 2~px.
With probability 0.10 a valid seed is moved 20 to 28~px in a random direction to give an example of a lost walk.
The target at each position comes from its projection onto the annotated path.
The state is lost if the position is more than 16~px from the path, and base if the remaining arc length is at most 4~px.
Otherwise the state is continue, and the target displacement points to the path point 8~px ahead, capped at 6~px.
The loss is the sum of an action term and a state term.
The action term is a smooth L1 loss~\cite{girshick2015fast} on the displacement plus one minus the cosine similarity of the directions.
The state term is a cross-entropy with class weights 2, 1, 2 and 2 for invalid, continue, base and lost.
Every 16 steps each walk may switch from the annotated displacement to its own predicted displacement, in the manner of scheduled sampling~\cite{bengio2015scheduled}.
The switch has probability $p_c = 0.30\,\min((c-1)/11, 1)$ in cycle $c$.
Two settings differ between training and decoding.
A walk has at most 96 steps in training and 128 at decoding.
The seed embedding is read at the rounded seed pixel in training and averaged over a radius of 2~px at decoding.

At decoding all seeds of a crop advance in parallel, and the state at each step is the one with the highest probability.
Table~\ref{tab:decoding} lists the settings for rejection, stopping and suppression.
At the first step a seed is rejected if the probability of the invalid state is at least the invalid-seed threshold, 0.6 for the published tracer.
If that probability is below the threshold but still the highest, the next most probable state is used.
An invalid decision at a later step counts as lost.
Arms are returned in the order of their seeds, in crop coordinates, and carry no anatomical number.

\begin{table}[!htbp]
\centering
\caption{\textbf{Decoding, stopping and suppression settings of the published tracer.}}
\label{tab:decoding}
{\small
\begin{tabular}{|>{\raggedright\arraybackslash}p{1.9in}|>{\raggedright\arraybackslash}p{4.1in}|}
\hline
\textbf{Setting} & \textbf{Value} \\ \thickhline
Seed & local maximum of the tip map with score $\geq$ 0.2; suppression radius 4~px \\ \hline
Seed embedding & mean over pixels with support $>$ 0.2 within 2~px of the seed \\ \hline
Step length & 0.5 to 6~px \\ \hline
Invalid-seed threshold & seed rejected at the first step if P(invalid) $\geq$ 0.6 \\ \hline
Stop: base & state ``terminate at base'' \\ \hline
Stop: lost & state ``terminate lost'' \\ \hline
Stop: border & position leaves the crop \\ \hline
Stop: length & path length $\geq$ 768~px \\ \hline
Stop: loop & position within 2~px of a path point more than 8 steps old \\ \hline
Stop: step limit & 128 steps \\ \hline
Returned arms & at least 2 points, not ended as lost, length $\geq$ 8~px \\ \hline
Smoothing & moving average over 5 points; both end points kept \\ \hline
Resampling & one point every 4~px, ordered from base to tip \\ \hline
Base score & score of the nearest peak of the base map (threshold 0.25, suppression radius 3~px) within 6~px of the path end; recorded, not used as a filter \\ \hline
Suppression: ranking & arms ranked by whether they reached a base, then by length, then by tip score \\ \hline
Suppression: duplicate & a lower-ranked arm is removed if its seed embedding is within 1.2 of a kept arm and at least 55\% of the vertices of the shorter arm lie within 4~px of the other arm \\ \hline
Suppression: fragment & a lower-ranked arm is removed if the embedding distance is $\leq$ 0.7, its length is $\leq$ 0.6 times that of the kept arm and the closest vertices are $\leq$ 12~px apart \\ \hline
\end{tabular}
}
\end{table}

\section*{Shortest-path baseline}

The shortest-path decoder reads the same maps as the tracer and uses no trained parameters.
Its nodes are the pixels with support above 0.5, and a search state is a pixel together with one of eight incoming directions.
Tips are peaks of the tip map (threshold 0.2, suppression radius 4~px) and bases are peaks of the base map (threshold 0.25, radius 3~px).
Peaks are snapped to the support mask within 8~px.
A step of length $l$ with unit direction $u$ costs
\begin{equation}
l\,\bigl(0.01 + 2.0\,(1-\bar{m}) + 0.5\,(1 + u\cdot t) + 2.0\,\max(\delta, 0) + 20.0\,\max(-\delta - 0.05, 0)\bigr) + 0.5\,(1 - u_{\mathrm{in}}\cdot u_{\mathrm{out}}) + 2.0\,l\,\lVert e(p) - e_{\mathrm{tip}}\rVert .
\label{eq:s1dijkstra}
\end{equation}
Here $\bar{m}$ is the mean centreline value of the two pixels, $t$ is the predicted tangent at the source pixel and $\delta$ is the change in progress per px.
The second term is the cost of a turn, and the third is the distance between the embedding of the destination pixel and the mean embedding within 2~px of the tip.
One Dijkstra search~\cite{dijkstra1959note} runs from each tip, using SciPy~\cite{virtanen2020scipy}.
Tips are paired with bases by the Hungarian method~\cite{kuhn1955hungarian,crouse2016implementing} on the path cost plus 2.0 times the embedding distance between tip and base.
A tip without a base is traced to the reachable pixel with the lowest progress.
A path is trimmed within 3~px of its base.
It is accepted if its length is at least 10~px, its mean step cost is at most 1.5 and its mean embedding distance to the tip is at most 1.5.
Its support component must also have at least 4\% of the area of the largest component.
Accepted paths are smoothed over 7 points, resampled every 4~px and passed through the suppression rules of Table~\ref{tab:decoding}.
All settings are fixed defaults and were not tuned.
One check on validation images, in which the tip threshold was raised from 0.2 to 0.4, increased the score of this decoder by about 1 percentage point.

\section*{Arm association in video}

In video, each sunflower sea star receives an identity from BoT-SORT~\cite{aharon2022botsort} as run by Ultralytics.
The settings are: high, low and new-track thresholds 0.25, 0.1 and 0.25; track buffer 30 frames; match threshold 0.8; score fusion on; camera-motion compensation by sparse optical flow; proximity threshold 0.5; appearance threshold 0.8; no re-identification model.
One arm tracker runs per star identity and is dropped after the star has not been seen for 30 frames.

The arm tracker links the arms of one star from frame to frame.
The body centre is the highest peak of the mouth map (threshold 0.15, suppression radius 8~px).
If there is none it is the centroid of the pixels with support above 0.3, and failing that the crop centre.
Each arm is expressed relative to the body centre and resampled to $K = 12$ points, evenly spaced along the arm from the tip.
For a traced arm $i$ and an existing track $j$, the shape distance $D_{ij}$ is the mean distance between corresponding points.
The points cover the arc length that both arms share, measured from the tip.
The cost of linking them is
\begin{equation}
D_{ij} = \frac{1}{K}\sum_{k=1}^{K}\bigl\lVert \tilde p_i(k)-\tilde p_j(k)\bigr\rVert_2,
\qquad
c_{ij} = D_{ij} + 2.0\,\lVert e_i - e_j\rVert_2 \quad \text{if } D_{ij}\leq 120\ \text{px},
\label{eq:s1assoc}
\end{equation}
where $e_i$ is the seed embedding of the arm and $e_j$ the stored embedding of the track.
Pairs with $D_{ij}$ above 120~px cannot be linked.
The Hungarian method assigns arms to tracks.
Each arm and each track can stay unmatched at a cost of 45, so a pair is accepted only if $c_{ij} < 90$.
A matched track takes the new polyline.
Its embedding is updated to 0.6 times the old plus 0.4 times the new embedding if the two differ by at most 1.0.
The tracker has no motion model: a track is compared at its last observed shape.
Unmatched arms start new tracks, and a track is deleted after more than 16 frames without a match.
No annotation of arm identity over time exists, so the accuracy of these links was not measured against ground truth.

\section*{Scoring and statistics}

\subsection*{Arms}

Annotated polylines are smoothed with two Chaikin iterations and resampled at 1~px, and traced polylines are resampled at 2~px.
The distance between a traced arm $P$ and an annotated arm $G$ is the symmetric Chamfer distance~\cite{barrow1977parametric,fan2017point},
\begin{equation}
d_{\mathrm{C}}(P,G) = \frac{1}{2}\left(
  \frac{1}{|P|}\sum_{p\in P}\min_{g\in G}\lVert p-g\rVert_2
  +
  \frac{1}{|G|}\sum_{g\in G}\min_{p\in P}\lVert p-g\rVert_2 \right).
\label{eq:s1chamfer}
\end{equation}
All pairs of a traced and an annotated arm in an image are sorted by this distance.
Pairs are accepted in ascending order if the distance is at most 8~px and neither arm is already matched.
Accepted pairs are true positives (TP), unmatched traced arms are false positives (FP) and unmatched annotated arms are false negatives (FN).
From these counts,
\begin{equation}
\mathrm{precision} = \frac{\mathrm{TP}}{\mathrm{TP}+\mathrm{FP}}, \qquad
\mathrm{recall} = \frac{\mathrm{TP}}{\mathrm{TP}+\mathrm{FN}}, \qquad
F_1 = \frac{2\,\mathrm{TP}}{\max\left(2\,\mathrm{TP}+\mathrm{FP}+\mathrm{FN},\ 1\right)} .
\label{eq:s1f1}
\end{equation}
Pooled F1 uses the counts summed over all images of a cohort.
Group-macro F1 sums the counts within each acquisition group, computes F1 per group and takes the unweighted mean over groups.
The score of a configuration is the mean of its group-macro F1 over three seeds.
The primary cohort is the development test set (493 images, 180 groups).

Configurations are compared with a paired bootstrap over acquisition groups~\cite{efron1979bootstrap,field2007bootstrapping}.
For each group we take the difference between two configurations in F1, averaged over seeds.
We draw the 180 groups with replacement 2,000 times (seed 20261026) and use the same draws for every contrast.
The 95\% interval runs from the 2.5th to the 97.5th percentile of the mean difference.
The two-sided p value centres the bootstrap distribution on zero,
\begin{equation}
p = \frac{1 + \#\{b : |\delta^{*}_b - \hat\delta| \geq |\hat\delta|\}}{B+1},
\label{eq:s1pvalue}
\end{equation}
with $\hat\delta$ the observed difference, $\delta^{*}_b$ the difference in draw $b$ and $B = 2{,}000$.
The p values of the eight registered contrasts are adjusted by the Holm method~\cite{holm1979simple}.
With 2,000 draws the smallest possible p value is 0.0005, so the smallest adjusted value over eight contrasts is 0.004.

\subsection*{Detections}

Predictions of the sunflower sea star class with a confidence of at least 0.25 are taken in descending order of confidence.
Each takes the unused annotated mask with the highest intersection over union and is a true positive if that value is at least 0.5.
Precision, recall and F1 follow Eq~\ref{eq:s1f1}.
Intervals for the detector resample capture groups 2,000 times (seed 20261023) and take percentiles.
The p value of a detector contrast with bootstrap differences $d_b$ is $\min\bigl(1,\ 2\,(\min(\#\{d_b \leq 0\}, \#\{d_b \geq 0\}) + 1)/(B+1)\bigr)$.
This differs from Eq~\ref{eq:s1pvalue}.
Seven contrasts on mask AP were registered for the detector and adjusted by the Holm method.

\section*{Detector results}

Table~\ref{tab:detdomains} gives the scores of the published detector on the test split.
On all 1,249 test images precision was 0.920, recall 0.903 and F1 0.911.
The mean absolute error of the star count per image was 0.138, with a bias of $-0.021$.
The detector made a false positive on 1 of the 550 test images without a sunflower sea star and on none of the 8 images without any annotation.

Recall on field images without PWS was 0.668 (123 of 184 stars in 7 capture groups).
Two sources account for most of the missed stars.
In the survey video (\texttt{rov\_svmp}; 32 images) the detector found 2 of 30 stars. \gap{platform and origin of the survey video}
In four images of a captive tank outside the laboratory (\texttt{field\_porch}) it found 51 of 80 stars.
On PWS all 243 stars were found, but this is a within-cruise result (see Datasets).
The Web row rests on 28 stars.
The campaign recorded bootstrap intervals for mask AP only, so Table~\ref{tab:detdomains} gives no interval for precision, recall or F1.

\begin{table}[!htbp]
\centering
\caption{\textbf{Published detector on the test split, by evaluation domain.} Size L, label scheme B, seed 0; confidence at least 0.25; mask intersection over union at least 0.5.}
\label{tab:detdomains}
{\footnotesize
\setlength{\tabcolsep}{3.5pt}
\begin{tabular}{|l|r|r|r|r|r|r|r|r|r|l|}
\hline
\textbf{Test images} & \textbf{Images} & \textbf{Groups} & \textbf{Stars} & \textbf{TP} & \textbf{FP} & \textbf{FN} & \textbf{Precision} & \textbf{Recall} & \textbf{F1} & \textbf{Mask AP (95\% CI)} \\ \thickhline
All & 1,249 & 520 & 1,413 & 1,276 & 111 & 137 & 0.920 & 0.903 & 0.911 & 0.719 (0.625 to 0.825) \\ \hline
Lab & 352 & 25 & 958 & 885 & 92 & 73 & 0.906 & 0.924 & 0.915 & 0.697 (0.637 to 0.784) \\ \hline
Field without PWS & 152 & 7 & 184 & 123 & 11 & 61 & 0.918 & 0.668 & 0.774 & 0.529 (none recorded) \\ \hline
PWS & 243 & 1 & 243 & 243 & 6 & 0 & 0.976 & 1.000 & 0.988 & 0.938 (one group) \\ \hline
Web & 502 & 487 & 28 & 25 & 2 & 3 & 0.926 & 0.893 & 0.909 & 0.819 (0.736 to 0.904) \\ \hline
\end{tabular}
}
\begin{flushleft}\small
Groups are capture groups.
Mask AP is the registered selection metric and is shown for reference; its intervals come from 2,000 bootstrap draws of capture groups.
PWS shares its cruise with 635 training images.
\end{flushleft}
\end{table}

The selection rule chose size L with scheme B by a small margin.
Its mean validation mask AP exceeded that of the runner-up by 0.0010 (95\% CI $-0.0043$ to 0.0110), and it ranked first in 38\% of the bootstrap draws.
On the test split, size L with scheme C scored 0.71709 and size L with scheme B 0.71707, and the six best configurations lay within 0.004.
The rank correlation between validation and test mask AP over the 20 configurations was 0.838 (Spearman).
We therefore make no claim that the selected size or scheme is better than its neighbours.
Of the seven registered contrasts, the steps from size N to S and from S to M under scheme F had Holm-adjusted p values of 0.007 each.
The other five contrasts, including the three between scheme F and the other schemes at size L, had adjusted p values of 0.585 or more.

Table~\ref{tab:schemes} compares the four label schemes at size L.
Recall on all test images was between 0.890 and 0.908.
On field images without PWS, recall was 0.669 to 0.672 for the three schemes that label other objects and 0.545 for scheme P.
This difference was not among the registered contrasts and is descriptive.

\begin{table}[!htbp]
\centering
\caption{\textbf{Four label schemes at detector size L.} Means over three seeds on the test split.}
\label{tab:schemes}
{\small
\begin{tabular}{|l|r|r|r|r|r|}
\hline
 & & \multicolumn{3}{c|}{\textbf{All test images}} & \textbf{Field without PWS} \\ \hline
\textbf{Scheme} & \textbf{Classes} & \textbf{Precision} & \textbf{Recall} & \textbf{F1} & \textbf{Recall} \\ \thickhline
F, every class separate & 37 & 0.942 & 0.905 & 0.923 & 0.669 \\ \hline
C, four coarse classes & 4 & 0.929 & 0.908 & 0.918 & 0.672 \\ \hline
B, sunflower sea star and other & 2 & 0.925 & 0.904 & 0.915 & 0.669 \\ \hline
P, sunflower sea star only & 1 & 0.940 & 0.890 & 0.914 & 0.545 \\ \hline
\end{tabular}
}
\begin{flushleft}\small
Detections at a confidence of at least 0.25, matched at a mask intersection over union of at least 0.5.
The published detector is seed 0 of scheme B.
\end{flushleft}
\end{table}

\section*{Tracer results}

\subsection*{Configurations and contrasts}

The tracer experiment comprised 27 trained components: nine kinds by three seeds.
These were three field networks, three field networks with a refined endpoint head and 21 tracers in seven configurations.
Every tracer was trained on the field network of the same seed.
The shortest-path decoder needs no training.
The invalid-seed threshold of each tracer configuration was chosen on validation images (see below): 0.6 for all configurations except the no-embedding and Sleepy Stars only configurations, which use none.

Table~\ref{tab:contrasts} gives the eight registered contrasts on the development test set.
The full tracer reached a group-macro F1 of 0.755, with precision 0.818, recall 0.773 and pooled F1 0.795 (means over three seeds).
The shortest-path decoder reached 0.705, with precision 0.725, recall 0.751 and pooled F1 0.737.
The difference was 5.01 percentage points (95\% CI 4.05 to 5.99; Holm-adjusted p = 0.004).

The no-embedding configuration changes two things at once.
It sets the seed embedding, the embedding channels and the similarity map to zero in the tracer input, and it also switches off duplicate and fragment suppression at decoding.
Its registered contrast ($-2.49$ percentage points) is therefore a combined effect.
We separated the two parts after the test split had been read, by decoding again on a CPU with each part switched alone.
These post hoc contrasts are marked in Table~\ref{tab:contrasts}.
Removing suppression alone lowered the score by 2.14 percentage points, and removing the embedding input alone by 0.60.
The CPU decoding differed from the registered readout by one false positive on 5 of 3,318 image rows, and gave $-2.51$ percentage points for the combined effect.
At a fixed set of walks, suppression removed 10.17\% of the traced arms of the full tracer (15,442 to 13,872 over three seeds).
It removed 150 true positives and 1,420 false positives, which changed precision from 0.744 to 0.818 and recall from 0.784 to 0.773.

\begin{table}[!htbp]
\centering
\caption{\textbf{Tracer configurations and contrasts on the development test set.} Group-macro F1 over 180 acquisition groups, mean over three seeds (SD over seeds). Differences are in percentage points.}
\label{tab:contrasts}
{\footnotesize
\setlength{\tabcolsep}{3.5pt}
\begin{tabular}{|>{\raggedright\arraybackslash}p{2.55in}|l|r|l|l|}
\hline
\textbf{Configuration and contrast} & \textbf{F1 (SD)} & \textbf{Difference} & \textbf{95\% CI} & \textbf{Holm p} \\ \thickhline
Full tracer, hidden width 128 (reference) & 0.7552 (0.0029) & & & \\ \thickhline
\multicolumn{5}{|l|}{\textbf{Registered contrasts}} \\ \hline
Shortest-path decoder $-$ full & 0.7051 (0.0061) & $-5.01$ & $-5.99$ to $-4.05$ & 0.004 \\ \hline
No embedding, input and suppression both off $-$ full & 0.7303 (0.0034) & $-2.49$ & $-3.04$ to $-1.89$ & 0.004 \\ \hline
No image features $-$ full & 0.7552 (0.0029) & $-0.005$ & $-0.39$ to 0.36 & 1.0 \\ \hline
No scheduled rollout in training $-$ full & 0.7576 (0.0030) & 0.23 & $-0.21$ to 0.68 & 1.0 \\ \hline
Hidden state reset at every step $-$ full & 0.7548 (0.0042) & $-0.04$ & $-0.45$ to 0.34 & 1.0 \\ \hline
Trained on the 403 Sleepy Stars crops only $-$ full & 0.7505 (0.0029) & $-0.47$ & $-0.97$ to 0.01 & 0.37 \\ \hline
Hidden width 64 $-$ full & 0.7562 (0.0018) & 0.10 & $-0.26$ to 0.45 & 1.0 \\ \hline
Full $-$ full on a field network with refined endpoint head & 0.7550 (0.0027) & 0.03 & $-0.09$ to 0.15 & 1.0 \\ \thickhline
\multicolumn{5}{|l|}{\textbf{Post hoc contrasts (CPU decoding, Holm over three tests)}} \\ \hline
Embedding input off, suppression on $-$ full & & $-0.60$ & $-1.09$ to $-0.14$ & 0.031 \\ \hline
Suppression off, embedding input on $-$ full & & $-2.14$ & $-2.53$ to $-1.75$ & 0.0015 \\ \hline
\end{tabular}
}
\begin{flushleft}\small
The F1 column gives the score of the configuration named first in each row; in the last registered row it is the score with the refined endpoint head.
F1 is given to four decimals so that small differences remain visible.
Intervals and p values come from 2,000 paired bootstrap draws of acquisition groups; registered p values are adjusted over eight contrasts.
For the endpoint refinement only the endpoint head of the field network was trained further, on the 403 Sleepy Stars training crops for 8 epochs.
The third post hoc test, embedding input off with suppression off in both arms, gave $-0.37$ ($-0.80$ to 0.07; Holm p = 0.098).
\end{flushleft}
\end{table}

\subsection*{Published tracer}

The published tracer is seed 0 of the full configuration with an invalid-seed threshold of 0.6.
Table~\ref{tab:cohorts} gives its scores by cohort and by the number of annotated arms.
On all 553 test images it matched 4,534 of 5,685 annotated arms (precision 0.821, recall 0.798, pooled F1 0.809).
On StarDome and StarTrack, which were not scored during development, pooled F1 was 0.915 and 0.868.
Each of these two cohorts has 30 images, and their test animals may also appear in training.
The campaign recorded one interval for a single cohort: the group-macro F1 on Sleepy Stars had a 95\% interval from 0.819 to 0.919 over five groups.

Pooled F1 was lowest for stars with 3 to 7 annotated arms (0.648, precision 0.605) and highest for stars with 14 to 16 arms (0.925).
The number of traced arms equalled the annotated number in 32.7\% of the images and was within one in 66.7\%.
The mean count error was $-0.29$ arms (median 0, mean absolute error 1.35).
The 4,534 matched arms lay a median of 2.74~px from the annotation by Chamfer distance (interquartile range 1.89 to 4.07~px), and 74.0\% lay within 4~px.
A matched arm is within 8~px by construction, so these values describe matched arms only.
The median error was 3.24~px at the tip (interquartile range 1.88 to 5.87~px) and 7.28~px at the base (4.32 to 11.53~px).

\begin{table}[!htbp]
\centering
\caption{\textbf{Published tracer on the test split, by cohort and by number of annotated arms.} Seed 0, invalid-seed threshold 0.6; arms matched at a Chamfer distance of at most 8~px.}
\label{tab:cohorts}
{\footnotesize
\setlength{\tabcolsep}{2.8pt}
\begin{tabular}{|l|r|r|r|r|r|r|r|r|r|r|}
\hline
\textbf{Test images} & \textbf{Images} & \textbf{Groups} & \textbf{Arms} & \textbf{TP} & \textbf{FP} & \textbf{FN} & \textbf{Precision} & \textbf{Recall} & \textbf{Pooled F1} & \textbf{Group-macro F1} \\ \thickhline
\multicolumn{11}{|l|}{\textbf{Cohort}} \\ \hline
All test images & 553 & 211 & 5,685 & 4,534 & 989 & 1,151 & 0.821 & 0.798 & 0.809 & 0.7778 \\ \hline
Development test set & 493 & 180 & 4,890 & 3,834 & 916 & 1,056 & 0.807 & 0.784 & 0.795 & 0.7554 \\ \hline
StarDome & 30 & 30 & 419 & 379 & 30 & 40 & 0.927 & 0.905 & 0.915 & 0.9092 \\ \hline
StarTrack & 30 & 1 & 376 & 321 & 43 & 55 & 0.882 & 0.854 & 0.868 & 0.8676 \\ \hline
Sleepy Stars & 63 & 5 & 650 & 560 & 59 & 90 & 0.905 & 0.862 & 0.883 & 0.8695 \\ \thickhline
\multicolumn{11}{|l|}{\textbf{Annotated arms per image, all test images}} \\ \hline
3 to 7 & 147 & & 878 & 613 & 400 & 265 & 0.605 & 0.698 & 0.648 & \\ \hline
8 to 10 & 184 & & 1,654 & 1,298 & 279 & 356 & 0.823 & 0.785 & 0.803 & \\ \hline
11 to 13 & 112 & & 1,349 & 1,010 & 197 & 339 & 0.837 & 0.749 & 0.790 & \\ \hline
14 to 16 & 68 & & 1,020 & 919 & 47 & 101 & 0.951 & 0.901 & 0.925 & \\ \hline
17 to 22 & 42 & & 784 & 694 & 66 & 90 & 0.913 & 0.885 & 0.899 & \\ \hline
\end{tabular}
}
\begin{flushleft}\small
Groups are acquisition groups.
Sleepy Stars is part of the development test set.
StarTrack is one group, so its group-macro F1 equals its pooled F1.
Precision, recall and pooled F1 are computed from the counts in the row; group-macro F1 is given to four decimals as recorded.
Arms per bin are TP plus FN.
\end{flushleft}
\end{table}

\subsection*{Choice of the invalid-seed threshold}

The threshold was chosen on the 478 validation images from five candidates: none, 0.3, 0.4, 0.5 and 0.6.
The criterion was the validation group-macro F1, averaged over three seeds.
Candidates within 0.002 of the best counted as tied, and ties went to no threshold and then to the lowest threshold.
For the full tracer the five candidates scored 0.7511, 0.7390, 0.7444, 0.7514 and 0.7537.
The rule chose 0.6 (0.752, 0.749 and 0.760 for seeds 0, 1 and 2).
The value 0.6 is the upper edge of the grid, and larger values were not tested.
The same rule kept hidden width 128 (0.7537) over width 64 (0.7505).

On the validation images the published tracer, with the threshold of 0.6, had precision 0.820, recall 0.784, pooled F1 0.801 and group-macro F1 0.7518.
With no threshold the same weights had precision 0.836, recall 0.773 and pooled F1 0.803.
The value 0.8034 that is stored with the checkpoint and quoted in the documentation of the inference package is this pooled validation F1 with no threshold.
It is not the accuracy of the published configuration.

\subsection*{Field network}

No measurement of the field network maps was made on the test split.
The training log gives validation scores for the published weights at the last epoch (478 crops).
The centreline path F1 was 0.4381 at a tolerance of 2~px, 0.6858 at 4~px (precision 0.7201, recall 0.6562) and 0.8921 at 8~px (precision 0.8951, recall 0.8905).
The arm support, thresholded at 0.5, had precision 0.571, recall 0.849 and an intersection over union of 0.519.
